\documentclass[11pt]{article}
\usepackage[margin=1in]{geometry}
\usepackage{amsmath, amssymb, bm}
\usepackage{hyperref}

\newcommand{\R}{\mathbb{R}}
\newcommand{\SO}{\mathrm{SO}(3)}
\newcommand{\vee}{\mathrm{vee}}
\newcommand{\skew}[1]{[#1]_\times}
\newcommand{\der}[1]{\dot{#1}}
\newcommand{\T}{\!\top}

\title{Closed-Loop Quadrotor Model for Reachability Analysis}
\author{cp\_reach}
\date{}

\begin{document}
\maketitle

\section{Overview}

This document describes the continuous-time closed-loop quadrotor model
defined in \texttt{quadrotor\_closed\_loop.mo}. The model is structured
in three layers via Modelica inheritance:

\begin{enumerate}
  \item \textbf{RigidBody} --- 6-DOF kinematics using rotation matrices
        (15 states)
  \item \textbf{Quadrotor} --- adds thrust mapping and angular dynamics
        (+3 states = 18)
  \item \textbf{QuadrotorClosedLoop} --- adds a geometric SO(3)
        controller with velocity integrators (+3 states = 21)
\end{enumerate}

All dynamics are polynomial or bilinear in the states---there are no
trigonometric functions---making the model suitable for polytopic
Jacobian-based reachability analysis in cp\_reach.

%% ====================================================================
\section{Rigid Body Kinematics}
%% ====================================================================

The rigid body state is $(\bm{p},\, \bm{v},\, R)$ where
$\bm{p}, \bm{v} \in \R^3$ are world-frame (NED) position and velocity,
and $R \in \SO$ is the body-to-world rotation matrix.

\begin{align}
  \der{\bm{p}} &= \bm{v} \\[4pt]
  \der{\bm{v}} &= R\,\bm{a} + \bm{g}, \qquad
    \bm{g} = \begin{bmatrix} 0 \\ 0 \\ g \end{bmatrix} \\[4pt]
  \der{R} &= R\,\skew{\bm{\omega}}
\end{align}

where $\bm{a} = [a_x,\, a_y,\, a_z]^\T$ is the body-frame specific
force (excluding gravity), $\bm{\omega} = [\omega_x,\, \omega_y,\,
\omega_z]^\T$ is the body-frame angular velocity, and
$\skew{\bm{\omega}}$ is the skew-symmetric matrix:
\begin{equation}
  \skew{\bm{\omega}} =
  \begin{bmatrix}
    0 & -\omega_z & \omega_y \\
    \omega_z & 0 & -\omega_x \\
    -\omega_y & \omega_x & 0
  \end{bmatrix}.
\end{equation}

In NED convention, positive $p_z$ is downward and gravity appears as
$+g$ in the $v_z$ equation.

%% ====================================================================
\section{Quadrotor Plant}
%% ====================================================================

The quadrotor extends the rigid body with a thrust-to-force mapping and
Euler's rotational equations. The control inputs at this level are thrust
$T$ and body-frame moments $\bm{M} = [M_x,\, M_y,\, M_z]^\T$.

\subsection{Force Mapping}

Thrust acts along the body $-z$ axis (upward in NED):
\begin{equation}
  \bm{a} = \begin{bmatrix} 0 \\ 0 \\ -T/m \end{bmatrix}.
\end{equation}

\subsection{Angular Dynamics}

Euler's equations with diagonal inertia tensor
$I = \mathrm{diag}(I_{xx},\, I_{yy},\, I_{zz})$:
\begin{align}
  I_{xx}\,\der{\omega}_x &= M_x + (I_{yy} - I_{zz})\,\omega_y\,\omega_z \\
  I_{yy}\,\der{\omega}_y &= M_y + (I_{zz} - I_{xx})\,\omega_x\,\omega_z \\
  I_{zz}\,\der{\omega}_z &= M_z + (I_{xx} - I_{yy})\,\omega_x\,\omega_y
\end{align}

The cross-coupling terms $(\omega_i \omega_j)$ are bilinear in the
states.

%% ====================================================================
\section{Geometric Controller}
%% ====================================================================

The controller follows the geometric tracking approach of Lee, Leok, and
McClamroch\footnote{T.~Lee, M.~Leok, N.~H.~McClamroch, ``Geometric
Tracking Control of a Quadrotor UAV on SE(3),'' \emph{IEEE CDC}, 2010.},
with a cascaded structure:
\[
  (\bm{p},\bm{v}) \;\xrightarrow{\text{PID}}\;
  \bm{a}_\mathrm{des} \;\xrightarrow{\text{project}}\;
  T \;\qquad\qquad
  R \;\xrightarrow{\text{SO(3)}}\;
  \bm{e}_R \;\xrightarrow{\text{rate PD}}\;
  \bm{M}
\]

\subsection{Outer Loop: Position and Velocity}

Define the tracking errors:
\begin{equation}
  \bm{e}_p = \bm{p} - \bm{p}_\mathrm{ref}, \qquad
  \bm{e}_v = \bm{v} - \bm{v}_\mathrm{ref}.
\end{equation}

A velocity-error integrator state $\bm{\xi} \in \R^3$ provides integral
action:
\begin{equation}
  \der{\bm{\xi}} = \bm{e}_v.
\end{equation}

The desired world-frame acceleration is computed via PID:
\begin{equation}
  \bm{a}_\mathrm{des} =
    -K_p\,\bm{e}_p - K_v\,\bm{e}_v - K_i\,\bm{\xi}
\end{equation}
where $K_p$, $K_v$, $K_i$ are diagonal gain matrices.

\subsection{Thrust Computation}

The thrust is the projection of the desired force onto the current body
$z$-axis:
\begin{equation}
  T = T_\mathrm{ff}
      - m\,(\bm{a}_\mathrm{des} - \bm{g})^\T (R\,\bm{e}_3)
\end{equation}
where $R\,\bm{e}_3 = [R_{13},\; R_{23},\; R_{33}]^\T$ is the third
column of $R$ (body $z$-axis in world frame) and $T_\mathrm{ff}$ is a
feedforward thrust term.

\subsection{Attitude Error: Geometric SO(3)}

Rather than using Euler angles (which introduce trigonometric
singularities), the attitude error is computed directly on $\SO$ using
the \emph{vee map}.

Define the relative rotation:
\begin{equation}
  S = R_\mathrm{ref}^\T R.
\end{equation}

The attitude error vector is extracted from the skew-symmetric part of
$S$:
\begin{equation}
  \bm{e}_R = \tfrac{1}{2}\,\vee(S - S^\T)
\end{equation}
where the vee map extracts the vector from a skew-symmetric matrix:
\begin{equation}
  \vee\!\begin{pmatrix}
    0 & -c & b \\ c & 0 & -a \\ -b & a & 0
  \end{pmatrix}
  = \begin{pmatrix} a \\ b \\ c \end{pmatrix}.
\end{equation}

Explicitly, the three error components are:
\begin{align}
  e_{R,x} &= \tfrac{1}{2}(S_{32} - S_{23}) \\
  e_{R,y} &= \tfrac{1}{2}(S_{13} - S_{31}) \\
  e_{R,z} &= \tfrac{1}{2}(S_{21} - S_{12})
\end{align}

For small rotations, $\bm{e}_R$ approximates the rotation vector between
$R$ and $R_\mathrm{ref}$. The computation involves only products of
rotation matrix entries---no trigonometric functions.

\subsection{Rate Error with Disturbance Injection}

The angular rate error incorporates a sensor disturbance
$\bm{d} = [d_p,\, d_q,\, d_r]^\T$ representing an acoustic attack bias
on the gyroscope:
\begin{equation}
  \bm{e}_\omega = (\bm{\omega} + \bm{d}) - \bm{\omega}_\mathrm{ref}.
\end{equation}

The controller observes $\bm{\omega} + \bm{d}$ rather than the true
$\bm{\omega}$, so the disturbance directly corrupts the rate feedback.

\subsection{Moment Commands}

The moments combine attitude proportional, rate derivative, feedforward,
and gyroscopic compensation terms:
\begin{align}
  M_x &= M_{x,\mathrm{ff}}
         - I_{xx}(K_{r,x}\,e_{R,x} + K_{\omega,x}\,e_{\omega,x})
         + (I_{yy} - I_{zz})\,\omega_y\,\omega_z \\
  M_y &= M_{y,\mathrm{ff}}
         - I_{yy}(K_{r,y}\,e_{R,y} + K_{\omega,y}\,e_{\omega,y})
         + (I_{zz} - I_{xx})\,\omega_x\,\omega_z \\
  M_z &= M_{z,\mathrm{ff}}
         - I_{zz}(K_{r,z}\,e_{R,z} + K_{\omega,z}\,e_{\omega,z})
         + (I_{xx} - I_{yy})\,\omega_x\,\omega_y
\end{align}

The gyroscopic terms cancel the cross-coupling in Euler's equations at
the reference trajectory, improving tracking performance.

%% ====================================================================
\section{cp\_reach Interface}
%% ====================================================================

The model follows the cp\_reach convention for closed-loop Modelica
files:

\begin{center}
\begin{tabular}{lll}
  \textbf{Category} & \textbf{Variables} & \textbf{Count} \\
  \hline
  States & $p_x, p_y, p_z, v_x, v_y, v_z,$
           $R_{11}\ldots R_{33},$
           $\omega_x, \omega_y, \omega_z,$
           $\xi_x, \xi_y, \xi_z$ & 21 \\[2pt]
  Reference inputs & $\bm{p}_\mathrm{ref}, \bm{v}_\mathrm{ref},
           R_\mathrm{ref}$ (9 entries),
           $\bm{\omega}_\mathrm{ref},$
           $T_\mathrm{ff}, \bm{M}_\mathrm{ff}$ & 19 \\[2pt]
  Disturbance inputs & $d_p, d_q, d_r$ & 3 \\[2pt]
  Algebraic (errors) & $\bm{e}_p, \bm{e}_v, \bm{e}_R, \bm{e}_\omega,
           \bm{a}_\mathrm{des}, S$ (9 entries) & 21 \\[2pt]
  Parameters & $m, I_{xx}, I_{yy}, I_{zz}, g,$
               gains (16 total) & 21 \\
\end{tabular}
\end{center}

\subsection{Error Dynamics for Reachability}

cp\_reach computes error dynamics by taking the Jacobian of the
closed-loop right-hand side with respect to the state:
\begin{equation}
  \der{\bm{e}} = J(\bm{x})\,\bm{e} + B_d\,\bm{d}
\end{equation}
where $J(\bm{x}) = \frac{\partial \bm{f}}{\partial \bm{x}}$ is
evaluated along the reference trajectory. Because all dynamics are
polynomial in the states, $J$ is also polynomial, and cp\_reach can
bound it over a polytopic region using vertex evaluation.

\subsection{Polytopic Considerations}

For a system with $n$ bounded state variables appearing nonlinearly in
the Jacobian, the polytopic approach evaluates $J$ at $2^n$ vertices.
The rotation matrix entries $R_{ij}$ and angular rates $\omega_i$
contribute 12 nonlinear state variables, giving up to $2^{12} = 4096$
vertices. This is computationally tractable but represents the main
scaling challenge for this model.

\end{document}
